\documentclass[aspectratio=169]{beamer}
\usepackage{ifxetex}
\usepackage{fontspec}
\usepackage{xunicode}
\usepackage{xltxtra}
\usepackage{xecyr}
\usepackage{polyglossia}
\usepackage{multicol}
\usepackage{pdfpages} 


\usepackage{beamerskoltech} 
% \renewcommand{\logoname}{sklogo.png} % <- default value for logo 
% \renewcommand{\logoname}{} % <- if you want no logo

\usepackage{progressbar}
\colorlet{progressbarcolor}{skoltechgreen}
\usepackage[mode=out]{inoutclass}
\definecolor{inclasscolor}{RGB}{46, 228, 182}
\definecolor{outclasscolor}{RGB}{41, 0, 204}
\usepackage{xargs}
\usepackage{latexLectures}
\usepackage{minted}
\usetikzlibrary{intersections}
\usetikzlibrary{calc}


\title{\LaTeX:\\ \Large from dummy to \TeX nician}
\subtitle{Command creation. How \TeX\ works-2}
\author{Anton Lioznov}
\institute{Skoltech, \\Project Center of Omics Technologies and Advanced Mass Spectrometry}
\date{ISP 2025,\\ \textit{lesson 5}}
\usetikzlibrary{decorations.pathreplacing}
\def\hasError{\inclass{
\begin{tikzpicture}[remember picture,overlay,shift={(current page.north east)}]
\draw[fill, red, anchor=north east,xshift=-1cm,yshift=-0.7cm] circle [radius=0.06];
\end{tikzpicture}}%
}
\begin{document}


\frame{\titlepage}

\AtBeginSection[] 
{
  \begin{frame}{What we will know?}
  \tableofcontents[currentsection,hideallsubsections]
  \end{frame}
}
\AtBeginSubsection[]
{
  \begin{frame}{What we will know?}
  \tableofcontents[currentsubsection, hideothersubsections, sectionstyle=show/hide, subsectionstyle=show/shaded/hide]
  \end{frame}
}

{\supressprogressbartrue


\begin{frame}\frametitle{What we will know?}
\tableofcontents[hideallsubsections]
\end{frame}

\supressprogressbarfalse
}

% \begin{frame}
% % \begin{enumerate}
% %     \item Overview and basis
% %         \begin{tikzpicture}[remember picture, overlay]
% %             \draw[decoration={brace}, decorate] 
% %             ([yshift=0.5ex]pic cs:start) --
% %             ([yshift=-2ex]pic cs:end)
% %             node[midway, right=0.3cm] {User level};
% %         \end{tikzpicture}\tikz[remember picture]\coordinate(start);
% %     \item Document creation\tikz[remember picture]\coordinate(end);
% %     \item TikZ and Typography
% %     \item \TeX\ and Typography
% %     \item Command creation
% % \end{enumerate}
     
% %      \begin{itemize}
% %     \item[$\left.\begin{cases}
% %         \text{1.} & \text{First item in the list} \\
% %         \text{2.} & \text{Second item in the list}
% %     \end{cases}\right\}$] Text explaining items 1 and 2
    
% %     \item[$\left.\begin{cases}
% %         \text{3.} & \text{Third item in the list}
% %     \end{cases}\right\}$]
    
% %     \item[$\left.\begin{cases}
% %         \text{4.} & \text{Fourth item in the list} \\
% %         \text{5.} & \text{Fifth item in the list}
% %     \end{cases}\right\}$] Text explaining items 4 and 5
% % \end{itemize}

% % \begin{enumerate}
% %     \item \tikz\node(item1){Item 1};
% %     \item \tikz\node(item2){Item 2};
% %     \item \tikz\node(item3){Item 3};
% %     \item \tikz\node(item4){Item 4};
% %     \item \tikz\node(item5){Item 5};
% % \end{enumerate}

% % \begin{tikzpicture}[overlay, remember picture]
% %     % Bracket for items 1 and 2
% %     \draw[decorate, decoration={brace, amplitude=5pt}, thick]
% %         ([xshift=-5pt, yshift=3pt]item1.north west) -- ([xshift=-5pt, yshift=-3pt]item2.south west)
% %         node[midway, left=12pt] {Group A};

% %     % Bracket for item 3
% %     \draw[decorate, decoration={brace, amplitude=5pt}, thick]
% %         ([xshift=-5pt, yshift=3pt]item3.north west) -- ([xshift=-5pt, yshift=-3pt]item3.south west)
% %         node[midway, left=12pt] {Solo Item};

% %     % Bracket for items 4 and 5
% %     \draw[decorate, decoration={brace, amplitude=5pt}, thick]
% %         ([xshift=-5pt, yshift=3pt]item4.north west) -- ([xshift=-5pt, yshift=-3pt]item5.south west)
% %         node[midway, left=12pt] {Group B};
% % \end{tikzpicture}
% \end{frame}

\section{Technical agreements}

\begin{frame}{Agreements}{I}\relax
     {\Large inclass/outclass versions}
     \begin{itemize}
          \item two slightly different versions for class and home
          \item class version is more interactive and contains less information
          \inclass{\item \inclasshigh{} this line will be shown only in class} 
          \outclass{\item \outclasshigh{} this line will be shown only at home version}
     \end{itemize}
\end{frame}
\inclassframe{\begin{frame}{Frame for a class}
     
\end{frame}}
\outclassframe{\begin{frame}{Frame for home}
     
\end{frame}}

{\supressfootnotefalse
\begin{frame}[fragile]{Agreements}{II}\relax
\newcommand{\tikzmark}[1]{\tikz[overlay,remember picture] \node (#1) {};}

{ \Large Footnotes }
\begin{tikzpicture}[overlay,remember picture]
\draw[->,ultra thick] (0,0.1) to[out=0,in=45] (8, -1.5) to[out=225,in=90] (5,-5.5);
\end{tikzpicture}

\skfootnote{Like this}
\begin{itemize}
     \item For second reading
     \item Contain advanced usage of the command 
     \item Contain references to read more 
     \begin{itemize}
         \item to the exact chapter 
         \item (often) with the href to exact page  
     \end{itemize}
     \item Contains some comments
     \item Mostly for outclass version
\end{itemize}
\end{frame}

}

{\inclassmodetrue
\newcommand{\tikzmark}[1]{\tikz[overlay,remember picture] \node (#1) {};}

\begin{frame}{Agreements}{III}\relax
{ \Large Additional information -- ``magic'' \tikzmark{startM}} 

\begin{tikzpicture}[remember picture,overlay,shift={(current page.north east)}]
\node[anchor=north east,xshift=-0cm,yshift=-0cm](endM) {%
{\includegraphics[width=1cm]{images/magic}}%
};
\end{tikzpicture}%

\begin{tikzpicture}[overlay, remember picture]
    \draw[->,ultra thick] (startM) to[out=0, in=180] (endM);
\end{tikzpicture}

\begin{itemize}
     \item To have the full picture 
     \item Not to analyze or to puzzle out in class 
\end{itemize}
\end{frame}

}


\begin{frame}{\exFrame{Agreements}}{V}\relax


{ \Large Exercises } \begin{tikzpicture}[overlay]
% \draw[white] (0,0) -- (0, 0.6);
\draw[->,ultra thick] (0,0.1) to[out=0,in=-90](-2.1,3.15);
\end{tikzpicture}


\begin{itemize}
     \item To work in class 
\end{itemize}
\end{frame}
\begin{frame}{Special thanks to}\relax

     Our TAs:
     \begin{itemize}
         \item Peter Borisovets
         \item Pavel Kuzmin
         \item Anna Litvin
     \end{itemize}

\end{frame}
\supressfootnotefalse
\ifinclasstrue
\def\skfootnote#1{\myfootnote{{\color{white!70!black}#1}}}
\fi

%%%%%%%%%%%%%%%%%%%%%%%%%% MAIN SECTION %%%%%%%%%%%%%%%%%%%%%%%%%%%%%%%%%


\section{Simple command creation and style files creation} 
% pro and contr
\begin{frame}[t, fragile]{Pros and Cons}
\vspace{-3.5ex}
\cwpa{
    \small
    \begin{itemize}
        \item[$+$] When you have lots of equations\vspace{-1ex}
        \item[$+$] When you have a complex, but typical document \vspace{-1ex}
        \item[$+$] When you carry about device-independant view\vspace{-1ex}
        \item[$+$] When you don't want not care about the beauty, but want it\vspace{-1ex}
        \item[$+$] When you are care about the beauty wery much\vspace{-1ex}
        \item[$+$] When you hate all these humanitarian GUI programs
        % \item[$+$] When you love text files 
        % \item[``+''] When you want to impress your partner
        % yeah, I leave the commented item for you, source-code reader!
    \end{itemize}
}{
    \small
    \begin{itemize}
        \item[$-$] when you want to put something in attributary position\vspace{-1ex}
        \item[$-$] when you want to do something ``against the rules''\vspace{-1ex}
        \item[$-$] when you want to work with visual-based things (tables,...)\vspace{-1ex}
        \item[$-$] when you want to do something really simple \vspace{-1ex}
        \item[$-$] when you want to do something ``quick and dirty'' \vspace{-1ex}
    \end{itemize}
}[t]

\cprotect\skfootnote{
\verb|\item[$+$]|
}
\end{frame}

\begin{frame}{Science reseach about \LaTeX}\relax
     \begin{center}
     ``We show that LaTeX users were {\color{red}slower} than Word users <...> and produced {\color{red}more typesetting},<...>. LaTeX users, however, more often report {\color{green}enjoying using} their respective software.''          
     \end{center}

     
\skfootnote{\url{https://journals.plos.org/plosone/article?id=10.1371/journal.pone.0115069} (or in russian \url{https://habr.com/post/375109/})}
\end{frame}

\begin{frame}{Illustration}\relax
\begin{center}
     
     \begin{tikzpicture}[axes/.style=,important line/.style={very thick},]
        \begin{scope}[axes]
        \draw[->] (0,0) -- coordinate (x axis mid) (8,0);
    	\draw[->] (0,0) -- coordinate (y axis mid) (0,5);
    	\node[below] at (x axis mid) {document complexity and size};
    	\node[rotate=90, above] at (y axis mid) {effort and time consumption};
        \end{scope}
        \draw (0, 1)  .. controls (3,1) and (6,2) .. (8, 4) node[below=0.5cm]{\LaTeX};
        \draw[dashed] (0, 0.1) .. controls (3, 0.5) and (4,2) .. (5, 5) node[left]{MS Word};
     \end{tikzpicture}
\end{center}
     \skfootnote{based on \url{http://www.pinteric.com/miktex.html} picture\\ tikZ for plotting \url{http://www.texample.net/tikz/examples/line-plot-example/}}
\end{frame}


\begin{frame}{Common belief}\relax
    \begin{center}
        \begin{tikzpicture}
             \node[align=center] (0,0) {
             \huge \LaTeX\ is only for use\\ \huge  in academic area
             };
             \uncover<2>{\node[rotate=30, bottom color=red!50, top color=red!50] (0,0) {\Huge WRONG}};
        \end{tikzpicture}
         
    \end{center}
    \skfootnote{It was done with \tt TikZ}
\end{frame}

\begin{frame}[t]{The power of \LaTeX\ in it's templates and flexability!}\relax
\vspace{-1ex}
     Look at examples at:
     \begin{itemize}
         \item \url{https://www.latextemplates.com/}
         \item \url{https://tex.stackexchange.com/questions/158668/nice-scientific-pictures-show-off}
         \item \url{https://tex.stackexchange.com/questions/1319/showcase-of-beautiful-typography-done-in-tex-friends}
        \item $\cdots$
     \end{itemize}
     
     \vspace{-1ex}
\inclassFrag{  Please, look at these sites for two minutes  }[-1]
     
     \skfootnote{here are some more links \stExC{https://tex.stackexchange.com/questions/1362/latex-template-gallery}
     }
\end{frame}

\begin{frame}{Conclusion}\relax 

Now, in \the\year , using \LaTeX\ to write scientific articles with no math inside is more metter of joy, not productivity: MS Office took over lots of \LaTeX 's ideas.\\[2ex]

But \LaTeX\ becoming better too! because of packages, online tools and developing \LaTeX 3.

And for something as complex as this presentation you'll spend way too more time, trying to reproduce it with MS Office.
     
     \skfootnote{If you compile this presentation from source, you'll find that the year in the slide -- is the year of compilation :)}
\end{frame}


% view
\begin{frame}{What we have}
\skfootnote{\wikiC{https://en.wikibooks.org/wiki/LaTeX/Introduction}\\ \normalfont\url{http://www.texample.net/tikz/examples/tree/}}
\begin{center}
\begin{tikzpicture}[sibling distance=10em,
  every node/.style = {shape=rectangle, rounded corners,
    draw, align=center,
    top color=white, bottom color=skoltechgreen!20}]]
  \node {\TeX}
    child { node {\LaTeX}
      child { node {XeLaTeX} }
      child { node {LuaTeX} }
      } ;
\end{tikzpicture}

\end{center}
\end{frame}

\note[itemize]{
\item Start from \LaTeX
}

% LaTeX
\begin{frame}{Definitions}
\skfootnote{\normalfont\url{https://ru.wikipedia.org/wiki/LaTeX}}
\begin{columns}
\begin{column}{0.5\textwidth}
\begin{tikzpicture}[sibling distance=10em,
  every node/.style = {shape=rectangle, rounded corners,
    draw, align=center,
    top color=white, bottom color=skoltechgreen!20}]]
  \node {\TeX}
    child { node[text=red] {\LaTeX}
      child { node {XeLaTeX} }
      child { node {LuaTeX} }
      } ;
\end{tikzpicture}
\end{column}

\begin{column}{0.5\textwidth}
\small
{\csk \LaTeX} --- is the most popular set of macro-extensions (or macro package) of the computer typesetting system \TeX, which facilitates a typesetting of complex documents.

\end{column}
\end{columns}
\end{frame}

\note[itemize]{
\item What is macro 
\item macros = shortcuts 
}

% TeX
\begin{frame}{Definitions}
\skfootnote{\normalfont\url{https://en.wikipedia.org/wiki/TeX}}
\begin{columns}
\begin{column}{0.5\textwidth}
\begin{tikzpicture}[sibling distance=10em,
  every node/.style = {shape=rectangle, rounded corners,
    draw, align=center,
    top color=white, bottom color=skoltechgreen!20}]]
  \node[text=red] {\TeX}
    child { node {\LaTeX}
      child { node {XeLaTeX} }
      child { node {LuaTeX} }
      } ;
\end{tikzpicture}
\end{column}

\note{
how to pronaunce\\
Book The Art of Computer Programming \\ 
History 1978 (pdf=1993)
}

\begin{column}{0.5\textwidth}
\small
{\csk \TeX} --- is a typesetting system designed and mostly written by Donald Knuth. TeX was designed with two main goals in mind: to allow anybody to produce high-quality books using minimal effort, and to provide a system that would give exactly the same results on all computers, at any point in time

\end{column}
\end{columns}
\end{frame}

% LuaTeX and XeLaTeX
\begin{frame}{Definitions\preMagicPage}
\skfootnote{\normalfont\url{https://en.wikipedia.org/wiki/LuaTeX}, \url{https://en.wikipedia.org/wiki/XeTeX}}
\begin{columns}
\begin{column}{0.5\textwidth}
\begin{tikzpicture}[sibling distance=10em,
  every node/.style = {shape=rectangle, rounded corners,
    draw, align=center,
    top color=white, bottom color=skoltechgreen!20}]]
  \node {\TeX}
    child { node {\LaTeX}
      child { node[text=red] {XeLaTeX} }
      child { node[text=red] {LuaTeX} }
      } ;
\end{tikzpicture}
\end{column}

\begin{column}{0.5\textwidth}
\small
{\csk XeLaTeX} --- XeTeX is a \TeX typesetting engine using Unicode and supporting modern font technologies such as OpenType, Graphite and Apple Advanced Typography

{\csk LuaTeX} --- LuaTeX is a \TeX-based computer typesetting system which started as a version of pdfTeX with a Lua scripting engine embedded
\end{column}
\end{columns}
\end{frame}



\begin{frame}{Resourses\magicPage}%{Books etc}

\vspace{-3.0ex}
\begin{itemize}
    \item Knuth ``The \TeX Book'' (en, ru)
    \item L'vovsky ``Nabor i verstka v sisteme \LaTeX'' (ru)
    \item Lamport. ``\LaTeX. A Document Preparation System, User’s Guide and Reference Manual'' (en)
    \item Gratzer ``Math into \LaTeX'' (en) 
    \item Oetiker ``The Not So Short Introduction to \LaTeX'' (en, ru)
    \item \url{https://www.overleaf.com/learn}
    \item \url{https://www.latex-project.org/help/}
    \item \url{https://texfaq.org/}
\end{itemize}
\end{frame}

\note{
We will use papeeria\\
pointed about footnotes and colors!
}

\begin{frame}[fragile]{Resourses \magicPage}{Interesting links}
\cprotect\skfootnote{
\verb|\begin{description} \item[questions about \TeX]\end{description}|
}
\small
\begin{description}
    \item[questions about \TeX] \url{https://tex.stackexchange.com} \vspace{-1ex}
    \item[knowing a command of the symbol] \url{http://detexify.kirelabs.org/classify.html} \vspace{-1ex}
    \item[beauty of TikZ] \url{http://www.texample.net/tikz/examples/} \vspace{-1ex}
    \item[beauty of pictures] \url{https://tex.stackexchange.com/questions/158668/nice-scientific-pictures-show-off}\vspace{-1ex}
    \item[beauty of typesetting] \url{https://tex.stackexchange.com/questions/1319/showcase-of-beautiful-typography-done-in-tex-friends}\vspace{-1ex}
\end{description}
\end{frame}

\begin{frame}[fragile]{where to get}
\begin{enumerate}
    \item Online
    \begin{itemize}
        \item \url{http://papeeria.com}
        \item \url{https://overleaf.com}
    \end{itemize}
    \item Offline
    \begin{itemize}
        \item \LaTeX{} \url{https://www.latex-project.org/get/}
        \item package manager \verb'tlmgr'
    \end{itemize}
\end{enumerate}
     
\end{frame}







\subsection{Command creation}
\graphicspath{{sec02/images/}}
\lstset{inputpath=sec02/code}

\begin{frame}{WYSiWYG vs not-WYSiWYG editors}
\cwpa{
{\csk WYSiWYG} -- \textit{What You See is What You Get} editor\\~\\

\inclassFrag{Examples?}[0]
}{
Microsoft Word
\insImg{wysiwyg.png}
}
\end{frame}


\begin{frame}[fragile]{not-WYSiWYG}


\cprotect\skfootnote{``\verb|\begin{frame}[fragile]|'' to put code inside}


\inclassFrag{Examples?}[1]
\cprotect[mm]\cwpa{
{\small HTML }

\begin{lstlisting}[language=html]
<html>
  <head>
    <meta charset="utf-8">
  </head>
  <body>
    <h1>Header</h1>
    <i>Hello</i>,<br/> world!  <!-- comment -->
  </body>
</html>
\end{lstlisting}

}{
{\small \LaTeX}
\lstinputlisting{helloworld.tex}

}[t]

\end{frame}

\begin{frame}[fragile]{Commands}
    \lstinline[basicstyle=\tt\large]|\command[o1, o2]{n1, n2=value}[o3]{n3}|\par
    (o = optional argument, n = necessary argument)~\\[2ex]

\pause
    \skfootnote{\knuthc{ch 3}, \lvoc{ch 2.3, 2.7}}

    {\csk Command symbols}
    \verb|\$ \# \{ \} \^{} \& \_ \~{} \\|
    \vfil
    {\csk Command words}
    \verb|\sin \LaTeX \Rightarrow \qquad|
    
    {\csk Enviruments}
    \verb|\begin{frame}\end{frame}  \begin{equation}\end{equation}|
    
\end{frame}

\begin{frame}[fragile]{Document structure}{overview}
    \lstinputlisting{structureOfDocument.tex}
\end{frame}

\begin{frame}[fragile]{Document structure}{class files}

Class of the document responsable for the large-scale settings
    \cprotect\skfootnote{\lvoc{IV.1}, \lamc{2.2.2} \url{https://texblog.org/2013/02/13/latex-documentclass-options-illustrated/}\\ envirament \lstinline|tabbing|} 
{
\scriptsize
\begin{tabbing}
\lstinline|\documentclass|\hspace{-1ex} \= \lstinline|[10pt,|  \= \lstinline|onecolumn,|  \=  \lstinline|a4paper]|\hspace{-1ex} \= \lstinline|{article}| \kill
\> \> \> \> \lstinline|{beamer} %presentation, poster| \\
\> \> \> \> \lstinline|{report}| \\
\> \> \> \> \lstinline|{book}| \\
\> \> \> \> \lstinline|{standalone} %for one picture/equation| \\
\> \> \> \> \lstinline|{extarticle} %if you want 14pt font size| \\
\lstinline|\documentclass| \> \lstinline|[10pt,|  \> \lstinline|onecolumn,|  \>  \lstinline|a4paper]| \> \lstinline|{article}| \\ 
\> \lstinline|[12pt] %fontsize| \> \\ 
\> \> \lstinline|[twocolumns] %number of columns in document| \> \\ 
\> \> \> \lstinline|[a5paper] %paper size| \> \\ 
\end{tabbing}}
\cprotect\skfootnote{\slshape in this presentation \verb|[aspectratio=169]|}
\end{frame}
\note{
but there is no ``fixed set'' of classes. Can build our own 
}

%%TSK: предложить создать документ и поиграться с параметрами
\begin{frame}[fragile]{Document structure}{style files}
     Style files are responsable for settings and providing new commands 
     
     \vfill
     \hspace{-1ex}
     \lstinline[basicstyle=\tt\large]|\usepackage[optional]{necessary}{packagename}|
     \vfill
\end{frame}

\note{practicully everything is from package }


\inclassFrame{
\begin{frame}{This slide means it is ``in-class task time!''}\relax
\centering
     \LARGE This slide means it is ``in-class task time!''
\end{frame}
}
\cprotect\inclassFrame{
\begin{frame}[fragile]{Try to compile your first doc}\relax
     \lstinputlisting[basicstyle=\tt\small]{first.tex}
\end{frame}
}
\subsection{Style and class files}
\graphicspath{{sec03/images/}{sec03/code/}}
\lstset{inputpath=sec03/code}
% \knuthc  knuth the TeXBook
% \lvoc   Lvovsky
% \lamc  lamport latex 

\subsection{Math}

% \knuthc  knuth the TeXBook
% \lvoc   Lvovsky
% \lamc  lamport latex 
% \slshape different font for footnote
\graphicspath{{sec03/images/s1/}{sec03/code/s1/}}
\lstset{inputpath=sec03/code/s1/}

\begin{frame}[fragile]{Going $\to$ Math}\relax

\cprotect\skfootnote{
\vspace{-0.5ex}
\lvoc{I.3.29}, \lamc{3.1}, \knuthc{16}[137] \\ {\slshape 
also try \verb|\(x=y\)| and \verb|\[x=y\]| and environment \verb|equation*|}\\
\textt \verb|\|lstinputlisting[firstline=7, lastline=14,linerange=\{8-10\}]\{file.tex\}
}

\cprotect\twocolImg{
Math environments

\lstinputlisting[linerange={15-16,19-19,22-25}]{mdollar.tex}
}{mdollar}

\end{frame}

\begin{frame}{Going $\to$ Math}\relax
{\centering \Large
``Because mathematics is supposedly expensive.''\par} 

\hfill \textcopyright D.~Knuth``the \TeX Book''
     \skfootnote{\knuthc{16}[138]}
\end{frame}


\begin{frame}[fragile]{Indexes}\relax

\newcommand{\appendTline}[3]{
\makebox[10em]{#1\hfill} \hfill\hfill \makebox[5em]{\hfill #2\hfill} \hfill \makebox[5em]{\hfill #3\hfill} \hfill\hfill \strut
    \hrule
    \vspace{1ex}
}

\cprotect[mm]\appendTline{upper ind}{\lstinline[basicstyle=\tt\normalsize,]|$x^2$|}{ $x^2$}

\cprotect[mm]\appendTline{lower ind}{\lstinline[basicstyle=\tt\normalsize,]|$x_2$|}{ $x_2$}
\cprotect[mm]\appendTline{lower and upper ind}{\lstinline[basicstyle=\tt\normalsize,]|$x^4_2$|}{ $x^4_2$}
\cprotect\inclassFrag{
Try to put more then one letter to the index}
\cprotect[mm]\appendTline{more letters in ind}{\lstinline[basicstyle=\tt\normalsize,]|$x_{ij}$|}{ $x_{ij}$}
\cprotect\inclassFrag{

How can you create ${}^3_2 He$?}[4]
\cprotect[mm]\appendTline{empty block}{\lstinline[basicstyle=\tt\normalsize,]|${}^3_2He$|}{ ${}^3_2He$}
\cprotect[mm]\appendTline{index in index}{\lstinline[basicstyle=\tt\normalsize,]|$x^{4^2}$|}{ $x^{4^2}$}

\skfootnote{\knuthc{16}[139]\\You really don't want to know, what code generate these lines...}
\end{frame}

\begin{frame}[fragile]{Fractions and (square) root}\relax
    \begin{columns}
        \begin{column}{0.45\textwidth}
          \hfill \lstinline[basicstyle=\tt\normalsize]|$\frac{x+z^2}{y-1}$| 
        \end{column}
        \begin{column}{0.45\textwidth}
             $$\frac{x+z^2}{y-1}$$
        \end{column}
    \end{columns}
    \vphantom.
    \hrule
    \vphantom) 
    \begin{columns}
        \begin{column}{0.45\textwidth}
          \hfill \lstinline[basicstyle=\tt\normalsize]|$\sqrt{x}$| 
        \end{column}
        \begin{column}{0.45\textwidth}
             $$\sqrt{x}$$
        \end{column}
    \end{columns}
    \vphantom(
    \hrule
    \vphantom) 
    \begin{columns}
        \begin{column}{0.45\textwidth}
          \hfill \lstinline[basicstyle=\tt\normalsize]|$\sqrt[y]{x}$| 
        \end{column}
        \begin{column}{0.45\textwidth}
             $$\sqrt[y]{x}$$
        \end{column}
    \end{columns}
    
\cprotect\skfootnote{\lmanc{16.7.3}[159] \lmanc{16.7.5}[160] \overC{https://www.overleaf.com/learn/latex/Fractions_and_Binomials} \slshape sometimes, if you want use fraction into fraction it is better to use \verb|\cfrac{x}{y}|. See typeography section or \lvoc{II.4.1, II.5.1}
also $\sphericalangle$ \verb|\genfrac|
}
\end{frame}

\begin{frame}[fragile]{Brackets\preMagicPage}\relax

You can't write just
\vphantom. 
    \begin{columns}
        \begin{column}{0.70\textwidth}
          \hfill \lstinline[basicstyle=\tt\normalsize]"$$(\frac{x}{y})$$" 
        \end{column}
        \begin{column}{0.3\textwidth}
             $$(\frac{x}{y})$$
        \end{column}
    \end{columns}
    \vphantom(
    \incPause
use \ccol{\left} and \ccol{\right}
     \vphantom. 
    \begin{columns}
        \begin{column}{0.70\textwidth}
          \hfill \lstinline[basicstyle=\tt\normalsize]"$$\left(\frac{x}{y}\right)$$" 
        \end{column}
        \begin{column}{0.3\textwidth}
             $$\left(\frac{x}{y}\right)$$
        \end{column}
    \end{columns}
    \vphantom(
    \incPause
or even like 
\vphantom. 
    \begin{columns}
        \begin{column}{0.70\textwidth}
          \hfill \lstinline[basicstyle=\tt\normalsize]"$$\left.\frac{x}{y}\right|_a^b$$" 
        \end{column}
        \begin{column}{0.3\textwidth}
             $$\left.\frac{x}{y}\right|_a^b$$
        \end{column}
    \end{columns}
    \vphantom(
\end{frame}

\begin{frame}[fragile]{Brackets\magicPage}\relax

Sometimes you need to manually set the bracket size. Then use something like this 
\vphantom. 
    \begin{columns}
        \begin{column}{0.70\textwidth}
          \hfill \lstinline[basicstyle=\tt\normalsize]"$$(\bigl( \Bigl( \biggl( \Biggl($$" 
        \end{column}
        \begin{column}{0.3\textwidth}
             $$(\bigl( \Bigl( \biggl( \Biggl($$
        \end{column}
    \end{columns}
    \vphantom(
    
    \vphantom. 
    \begin{columns}
        \begin{column}{0.70\textwidth}
          \hfill \lstinline[basicstyle=\tt\normalsize]"$$]\bigr] \Bigr] \biggr] \Biggr]$$" 
        \end{column}
        \begin{column}{0.3\textwidth}
             $$]\bigr] \Bigr] \biggr] \Biggr]$$
        \end{column}
    \end{columns}
    \vphantom(

\skfootnote{\knuthc{17}[147] \overC{https://www.overleaf.com/learn/latex/Brackets_and_Parentheses}}
     
\end{frame}

\inclassFrame{
\begin{frame}{Type these equations}{simple}\relax
\begin{equation}
     x^2-2y = h_{l_0^k}\tag{1}
\end{equation}
\begin{equation}
     \sqrt{h/(2a)}=t\tag{2}
\end{equation}

\begin{equation}
     \frac{x+y}{z_2 - \frac{1}{2}}=m\tag{3}
\end{equation}

\begin{equation}
     \left(\frac{x+y}{z}\right)=q\tag{4}
\end{equation}

\end{frame}
}

\inclassFrame{
\begin{frame}{Type these equations}{harder}\relax
\begin{equation}
    T_{N_1}(t) = \left.\frac{d}{dx_i} f^{\sqrt{n}}(z)\right/\left(\sqrt{\frac{m+e}{a}}-N\right) \tag{1}
\end{equation}

\end{frame}
}

\begin{frame}[fragile]{Text inside equations}\relax
Sometimes, you need to write a text \textit{inside} an equation

$$\frac{x+1}{y} = z; \hbox{if $x - 1 < y$, but not always!}$$

\cprotect\inclassFrag{try to type \lstinline[basicstyle=\tt]|$$\frac{x+1}{y} = z; if x - 1 < y, but not always!$$|}
but the direct solution remove all spaces! 
$$\frac{x+1}{y} = z; if x - 1 < y, but not always!$$
\incPause
Use \ccol{\hbox}:

\lstinline[basicstyle=\tt]|$$\frac{x+1}{y} = z; \hbox{if $x - 1 < y$, but not always!}$$| \incPause \hfill or~\ccol{\text}~from~{\csk amsmath}~package.  

\end{frame}

\inclassFrame{
\begin{frame}{g}\relax
     How will you type $g = 9.8 \frac{\hbox{m}}{\hbox{s}^2}$?
\end{frame}
}

\begin{frame}[fragile]{More symbols\preMagicPage}\relax
    \line{\hss\Large \url{http://detexify.kirelabs.org/classify.html}\hss}
    
    \inclassFrag{\centering
    Try to find the following symbols: \\ $\exists$,  $\to$, $\notin$, $\sim$, $\equiv$
    }[-1]
    
    \skfootnote{\vspace{-3ex}
    \wikiC{https://en.wikibooks.org/wiki/LaTeX/Special_Characters} \\
    \url{ftp://sunsite.icm.edu.pl/pub/CTAN/info/symbols/comprehensive/symbols-a4.pdf} (338 pages) \\ \url{https://www.rpi.edu/dept/arc/training/latex/LaTeX_symbols.pdf} (22 pages)}
\end{frame}

\begin{frame}[fragile]{More symbols\magicPage}{Greek letters}\relax
\footnotesize

\newlength{\myboxlen}%
\setlength{\myboxlen}{9em}

\newcommand{\showmsym}[1]{%
\makebox[\myboxlen]{\hfill\makebox[0.45\myboxlen]{\csk\hfill\string#1}\makebox[0.25\myboxlen]{$\to$}\makebox[0.3\myboxlen]{$#1$\hfill}\hfill}%
}
\showmsym{\alpha}\hfill
\showmsym{\beta}\hfill
\showmsym{\gamma}\hfill
\showmsym{\delta}\hfill
\showmsym{\epsilon}\hfill
\showmsym{\zeta}\hfill
\showmsym{\eta}\hfill
\showmsym{\iota}\hfill
\showmsym{\kappa}\hfill
\showmsym{\lambda}\hfill
\showmsym{\mu}\hfill
\showmsym{\nu}\hfill
\showmsym{\xi}\hfill
\showmsym{\tau}\hfill
\showmsym{\upsilon}\hfill
\showmsym{\rho}\hfill
\showmsym{\sigma}\hfill
\showmsym{\pi}\hfill
\showmsym{\phi}\hfill
\showmsym{\chi}\hfill
\showmsym{\psi}\hfill
\showmsym{\omega}\hfill
\makebox[\myboxlen]{}\hfill
\makebox[\myboxlen]{}\hfill
\makebox[\myboxlen]{}\hfill

\showmsym{\varepsilon}\hfill
\showmsym{\vartheta}\hfill
\showmsym{\varkappa}\hfill
\showmsym{\varrho}\hfill
\showmsym{\varphi}\hfill

~\\
\showmsym{\Gamma}\hfill
\showmsym{\Delta}\hfill
\showmsym{\Lambda}\hfill
\showmsym{\Xi}\hfill
\showmsym{\Sigma}\hfill
\showmsym{\Upsilon}\hfill
\showmsym{\Pi}\hfill
\showmsym{\Phi}\hfill
\showmsym{\Psi}\hfill
\showmsym{\Omega}\hfill
\makebox[\myboxlen]{}\hfill
\makebox[\myboxlen]{}\hfill
\makebox[\myboxlen]{}\hfill

\end{frame}

\begin{frame}[fragile]{More symbols\magicPage}{Other common used symbols}\relax
\footnotesize

\newlength{\myboxlen}%
\setlength{\myboxlen}{9em}

\newcommand{\showmsym}[1]{%
\makebox[\myboxlen]{\hfill\makebox[0.45\myboxlen]{\csk\hfill\string#1}\makebox[0.25\myboxlen]{$\to$}\makebox[0.3\myboxlen]{$#1$\hfill}\hfill}%
}
\showmsym{\infty}\hfill
\showmsym{\nabla}\hfill
\showmsym{\forall}\hfill
\showmsym{\partial}\hfill
\showmsym{\backslash}\hfill
\showmsym{\aleph}\hfill
\showmsym{\hbar}\hfill
\showmsym{\Re}\hfill
\showmsym{\Im}\hfill
\showmsym{\cdot}\hfill
\showmsym{\cdots}\hfill
\showmsym{\ldots}\hfill
\showmsym{\to}\hfill
\showmsym{\times}\hfill
\makebox[\myboxlen]{}\hfill
\makebox[\myboxlen]{}\hfill
\makebox[\myboxlen]{}\hfill

\normalsize
 \showmsym{\imath}\hfill \showmsym{\jmath}\hfill are useful for accents.

\end{frame}

\begin{frame}[fragile]{More symbols\magicPage}{Accents}\relax
% \footnotesize

\skfootnote{\knuthc{17}}
\newlength{\myboxlen}%
\setlength{\myboxlen}{9em}

\newcommand{\showmacc}[2]{%
\makebox[\myboxlen]{\hfill\makebox[0.45\myboxlen]{\hfill{\csk\string#1}\{#2\}}\makebox[0.25\myboxlen]{$\to$}\makebox[0.3\myboxlen]{$#1{#2}$\hfill}\hfill}%
}
\showmacc{\hat}{a}\hfill
\showmacc{\check}{a}\hfill
\showmacc{\tilde}{a}\hfill
\showmacc{\acute}{a}\hfill
\showmacc{\grave}{a}\hfill
\showmacc{\dot}{a}\hfill
\showmacc{\ddot}{a}\hfill
\showmacc{\breve}{a}\hfill
\showmacc{\vec}{a}\hfill
~\\[2ex]
\pause

\setlength{\myboxlen}{12em}
\Large \showmacc{\check}{a}\hfill \to \hfill
\setlength{\myboxlen}{13em}
{\csk \textbackslash{}skew5\textbackslash{}check}\{a\}{$\to$}$\skew5\check{a}$%

\end{frame}


\begin{frame}[fragile]{More symbols\magicPage}{Math fonts}\relax
    \begin{tabbing}
    ааaaaaaaaaaaaaaaaaaaaaaaaaaaaa \= aaaaaaaaaaaaaaaaaaaaaaaaaaaaaa \= \kill
    \texttt{\textbackslash mathrm\{\textit{letters,etc}\}} \> $\mathrm{ABCabc,123,\hat{a},\tilde{b},\tilde{c},\Psi\Omega}$ \\
    \texttt{\textbackslash mathbf\{\textit{letters,etc}\}} \> $\mathbf{ABCabc,123,\hat{a},\tilde{b},\tilde{c},\Psi\Omega}$ \\
    \texttt{\textbackslash mathsf\{\textit{letters,etc}\}} \> $\mathsf{ABCabc,123,\hat{a},\tilde{b},\tilde{c},\Psi\Omega}$ \\
    \texttt{\textbackslash mathtt\{\textit{letters,etc}\}} \> $\mathtt{ABCabc,123,\hat{a},\tilde{b},\tilde{c},\Psi\Omega}$ \\
    \texttt{\textbackslash mathit\{\textit{letters,etc}\}} \> $\mathit{ABCabc,123,\hat{a},\tilde{b},\tilde{c},\Psi\Omega}$ \\
    \texttt{\textbackslash mathnormal\{\textit{letters,etc}\}} \> $\mathnormal{ABCabc,123,\hat{a},\tilde{b},\tilde{c},\Psi\Omega}$ \\
    \texttt{\textbackslash mathcal\{\textit{capital letters}\}} \> $\mathcal{ABC}$ \\ 
    \texttt{\textbackslash mathds\{\textit{capital letters}\}} \> $\mathds{ABCRN}$ 
    \end{tabbing}
    \cprotect\skfootnote{\lstinline|\usepackage{dsfont}| for the last}
\end{frame}

\begin{frame}[fragile]{More symbols\magicPage}{Limiters}\relax
\skfootnote{\knuthc{17}}
\newlength{\myboxlen}%
\setlength{\myboxlen}{6em}
% \small
\newcommand{\showmsym}[1]{%
\makebox[\myboxlen]{\hfill\makebox[0.45\myboxlen]{\small\csk\hfill\string#1}\makebox[0.25\myboxlen]{$\to$}\makebox[0.3\myboxlen]{$$#1$$\hfill}\hfill}%
}
\newcommand{\showidenti}[1]{%
\makebox[\myboxlen]{\hfill\makebox[0.45\myboxlen]{\csk\hfill$#1$}\makebox[0.25\myboxlen]{$\to$}\makebox[0.3\myboxlen]{$$#1$$\hfill}\hfill}%
}
{\bfseries Brackets}

\showidenti{(}\hfill
\showidenti{)}\hfill
\showmsym{\langle}\hfill
\showmsym{\rangle}\hfill
\showidenti{[}\hfill
\showidenti{]}\hfill
\showmsym{\lbrack}\hfill
\showmsym{\rbrack}\hfill
\showmsym{\{}\hfill
\showmsym{\}}\hfill
\showmsym{\lbrace}\hfill
\showmsym{\rbrace}\hfill
\showmsym{\lfloor}\hfill
\showmsym{\rfloor}\hfill
\showmsym{\lceil}\hfill
\showmsym{\rceil}\hfill

{\bfseries Other}
\renewcommand{\to}{}

\showidenti{/}\hfill
\showmsym{\backslash}\hfill
\showidenti{|}\hfill
\showmsym{\vert}\hfill
\showmsym{\|}\hfill
\showmsym{\Vert}\hfill
\showmsym{\updownarrow}\hfill
\showmsym{\Updownarrow}\hfill
\showmsym{\uparrow}\hfill
\showmsym{\Uparrow}\hfill
\showmsym{\downarrow}\hfill
\showmsym{\Downarrow}\hfill

\end{frame}

\begin{frame}[fragile]{More symbols\magicPage}{Limiters with {\csk \string\Big} prefix}\relax
\skfootnote{\knuthc{17}[147]}
\newlength{\myboxlen}%
\setlength{\myboxlen}{6em}
% \small
\newcommand{\showmsym}[1]{%
\makebox[\myboxlen]{\hfill\makebox[0.45\myboxlen]{\scriptsize\csk\hfill\string#1}\makebox[0.25\myboxlen]{$\to$}\makebox[0.3\myboxlen]{$$\Big#1$$\hfill}\hfill}%
}


\showmsym{\vert}\hfill
\showmsym{\Vert}\hfill
\showmsym{\arrowvert}\hfill
\showmsym{\Arrowvert}\hfill
\showmsym{\lgroup}\hfill
\showmsym{\rgroup}\hfill
\showmsym{\lmoustache}\hfill
\showmsym{\rmoustache}\hfill
\showmsym{\bracevert}\hfill

\end{frame}

\begin{frame}[fragile]{More symbols\magicPage}{Operators}\relax
\skfootnote{\overC{https://www.overleaf.com/learn/latex/Integrals,_sums_and_limits}}

\newlength{\myboxlen}%
\setlength{\myboxlen}{9em}

\newcommand{\showmsym}[1]{%
\makebox[\myboxlen]{\hfill\makebox[0.45\myboxlen]{\csk\hfill\string#1}\makebox[0.25\myboxlen]{$\to$}\makebox[0.3\myboxlen]{$$#1$$\hfill}\hfill}%
}
\showmsym{\prod}\hfill
\showmsym{\sum}\hfill
\showmsym{\int}\hfill
\showmsym{\iint}\hfill
\showmsym{\oint}\hfill
\showmsym{\idotsint}\hfill

\end{frame}

\begin{frame}[fragile]{More symbols}{Operators}\relax
Mathematical tradition: write not $sin (x)$, but {\csk$\sin (x)$}.\incPause

\begin{center}
     \Large use \ccol{\sin}
\end{center}
\incPause
sometimes more effect is displayed:
\begin{columns}

\begin{column}{0.5\textwidth}
\hfill\lstinline[basicstyle=\tt]|$$\min_{x\to 0} f$$|     
\end{column}
\begin{column}{0.5\textwidth}
$$\min_{x\to 0} f$$
\end{column}
\end{columns}


    \skfootnote{\overC{https://www.overleaf.com/learn/latex/Operators} \knuthc{18}[173] \url{https://oeis.org/wiki/List_of_LaTeX_mathematical_symbols}}
\end{frame}

\begin{frame}[fragile]{Multiline equations\preMagicPage}\relax
     \cprotect\twocolImg{
        \lstinputlisting[linerange={9-14}]{arraySimple}
    }{arraySimple}

\incPause
notice {\csk \{clcl\}}, {\csk \&}, \ccol{\\}
\inclassFrag{Does it reminds you something?}[3]

\ccol{array} is a \ccol{tabular} for math mode!

\cprotect\skfootnote{{ } \lvoc{II.3}[75], \lvoc{II.4.2}[84] \lmanc{8.2}[55]}
\end{frame}

\begin{frame}[fragile]{Also array\magicPage}\relax
     \cprotect\twocolImg{
        \lstinputlisting[linerange={11-17}]{matrixesno}
    }{matrixesno}
    \cprotect\twocolImg{
    \lstinputlisting[linerange={9-16}]{equcaseno}
    }{equcaseno}

\cprotect\skfootnote{advanced: \knuthc{18}[186] \tugC{https://www.tug.org/utilities/plain/cseq.html\#vcenter-rp}}
\end{frame}

\begin{frame}[fragile]{Multiline equations\preMagicPage}\relax
 it is better to use 
 \begin{center}
 \Large \csk \verb|\usepackage{amsmath}|
 \end{center}
 
\cprotect\skfootnote{\normalfont \url{http://mirrors.mi.ras.ru/CTAN/macros/latex/required/amsmath/amsmath.pdf} \url{http://texdoc.net/texmf-dist/doc/latex/amsmath/amsldoc.pdf}}
\end{frame}


\begin{frame}[fragile]{Formula in multiple line\magicPage}\relax
\skfootnote{{ }\lvoc{II.4.2}[81]\\ You can modify the line gaps, push string to left or right,..}

    \cprotect\twocolImg{
\lstinputlisting[linerange={13-14, 17-20}]{multilines.tex}
}{multilines}

\end{frame}


\begin{frame}[fragile]{Multiple formulas\magicPage}\relax

% \vspace{-10ex}
    \cprotect\twocolImg{
\lstinputlisting[linerange={13-14, 17-21}]{gathermy.tex}
}{gathermy}

\incPause
% \vspace{-10ex}
notice \lstinline[basicstyle=\tt\normalsize]|\notag| !

\cprotect\skfootnote{{} \lvoc{II.4.2}[81]\\ \slshape  also see \verb'gather*' to ban numeration}

\end{frame}

\begin{frame}[fragile]{Multiple formulas and lines: alignment\magicPage}\relax
\cprotect\twocolImg{
    \lstinputlisting[linerange={13-14, 17-20}]{alignmy.tex}
}{alignmy}
\cprotect\twocolImg{
    \lstinputlisting[linerange={13-14, 17-22}]{splitmy}
}{splitmy}


ampersand {\csk \&} is stands for indent (as in tables)
\cprotect\skfootnote{{ }\lvoc{II.4.2}[82]\\  also see \verb'align*' to ban numeration}
\end{frame}

\begin{frame}[fragile, t]{Text inside equations\magicPage}\relax

\cprotect\twocolImg{
    \lstinputlisting[linerange={13-14, 17-23}]{alignetextwrong1.tex}
}{alignetextwrong1}

Problem:
\cprotect\twocolImg{
    \lstinputlisting[linerange={13-14, 17-23}]{alignetextwrong.tex}
}{alignetextwrong}

\end{frame}
\begin{frame}[fragile]{Text inside equations\magicPage}{solution}\relax
% \vspace{-1.3cm}
\cprotect\twocolImg{
    \lstinputlisting[linerange={13-14, 17-23}]{alignetext.tex}
}{alignetext}
\ccol{\text} and \ccol{\intertext}
% \inclassFrag{How will you type $g = 9.8 \frac{\text{m}}{\text{s}^2}$?}[-1]
\skfootnote{{} \lvoc{II.4.2}[83]}
\end{frame}

\begin{frame}[fragile]{System of equations\magicPage}\relax
% \vspace{-1.3cm}
\cprotect\twocolImg{
    \lstinputlisting[linerange={13-14, 17-24}]{equsis.tex}
}{equsis}
% \vspace{-2cm}
\cprotect\twocolImg{
    \lstinputlisting[linerange={13-14, 17-23}]{equcase}
}{equcase}
\skfootnote{{} \lvoc{II.4.2}[84]}
\end{frame}


\begin{frame}[fragile]{Matrix\magicPage}\relax
    \cprotect\twocolImg{
    \lstinputlisting[linerange={7-7, 11-18}]{matrixes}
}{matrixes}
\cprotect\skfootnote{{ } \lvoc{II.3}[74]\\ \slshape Also see \lstinline!matrix! without borders and \lstinline!vmpatrix! for |border|}
     
\end{frame}

\cprotect\inclassFrame{
\begin{frame}[fragile]{Try to type the following}\relax
\small
\verb!$$\sum_a^b x$$!, \hfill\verb!$$\sum\nolimits_a^b x$$!, \hfill\verb!$$\sum\limits_a^b x$$!

\verb!$\sum_a^b x$!, \hfill\verb!$\sum\nolimits_a^b x$!, \hfill\verb!$\sum\limits_a^b x$!

\verb!$$\int_a^b x$$!, \hfill\verb!$$\int\nolimits_a^b x$$!, \hfill\verb!$$\int\limits_a^b x$$!

\verb!$\int_a^b x$!, \hfill\verb!$\int\nolimits_a^b x$!, \hfill\verb!$\int\limits_a^b x$!
~\\[2ex]
compare the index position and the operator size
     
\end{frame}
}

\begin{frame}[fragile]{One over another\preMagicPage}{operators}\relax

\newcommand{\appendTline}[2]{\vspace*{10pt}\begin{columns}
        \begin{column}{0.45\textwidth}
          \hfill #1 
        \end{column}
        \begin{column}{0.45\textwidth}
             \hfill #2\hfill \hfill
        \end{column}
    \end{columns}
    \vphantom.
    \hrule
    }

    \cprotect[mm]\appendTline{\csk source}{\csk operator}
    \hrule height 0.05pt
    \cprotect[mm]\appendTline{\lstinline[basicstyle=\tt\normalsize,]|$$\int\limits_0^\pi$$|}{$$\int\limits_0^\pi$$}
    \cprotect[mm]\appendTline{\lstinline[basicstyle=\tt\normalsize,]|$$\int\nolimits_0^\pi$$|}{$$\int\nolimits_0^\pi$$}

\end{frame}

\begin{frame}[fragile]{One over another\magicPage}\relax

\newcommand{\appendTline}[2]{\vspace*{10pt}\begin{columns}
        \begin{column}{0.45\textwidth}
          \hfill #1 
        \end{column}
        \begin{column}{0.45\textwidth}
             \hfill #2\hfill \hfill
        \end{column}
    \end{columns}
    \vphantom.
    \hrule
    }

    \cprotect[mm]\appendTline{\csk source}{\csk result}
    \hrule height 0.05pt
    \cprotect[mm]\appendTline{\lstinline|$\stackrel {\Leftrightarrow}{A}$|}{$\stackrel {\Leftrightarrow}{A}$}
    \cprotect[mm]\appendTline{\lstinline|$A \stackrel{a'}{\rightarrow} D$|}{$A \stackrel{a'}{\rightarrow} D$}
    \cprotect[mm]\appendTline{\lstinline|$$\sum_{\substack{i\in[0;n]\\j\in[0;m]}} a_{ij}$$|}{$$\sum_{\substack{i\in[0;n]\\j\in[0;m]}} a_{ij}$$}
    
\end{frame}

\begin{frame}[fragile]{One over another\magicPage}\relax

\newcommand{\appendTline}[2]{\vspace*{10pt}\begin{columns}
        \begin{column}{0.45\textwidth}
          \hfill #1 
        \end{column}
        \begin{column}{0.45\textwidth}
             \hfill #2\hfill \hfill
        \end{column}
    \end{columns}
    \vphantom.
    \hrule
    }

    \cprotect[mm]\appendTline{\csk source}{\csk result}
    \hrule height 0.05pt
    
    \cprotect[mm]\appendTline{\lstinline|$\underbrace{a+\overbrace{b+c}+d}_{m}$|}{$\underbrace{a+\overbrace{b+c}+d}_{m}$}
    
    \cprotect[mm]\appendTline{\lstinline|$\lefteqn{\overbrace{ \phantom{1+2+3}}} 1+\underbrace{2+3+4}$|}{$\lefteqn{\overbrace{ \phantom{1+2+3}}} 1+\underbrace{2+3+4}$}
    
\end{frame}

\begin{frame}{Domain-specific packages}{Lots of them!}\relax
You can use 
    \begin{description}
        \item[Physics] \url{https://ctan.org/pkg/physics}
        \item[Chemistry] \url{http://www.mychemistry.eu/known-packages/}, \url{https://ru.overleaf.com/learn/latex/Chemistry_formulae}, \url{https://ctan.org/pkg/mhchem}
        \item[Biology] \url{https://www.tug.org/pracjourn/2007-4/senthil/senthil.pdf} 
    \end{description}
     
\end{frame}


\subsection{Text}

\graphicspath{{sec02/images/}{sec02/code/}}
\lstset{inputpath=sec02/code/}

\begin{frame}[label=style,fragile]{Built--in themes}
    \visible<1>{\hyperlink{simple}{\beamerbutton{SIMPLE}} \\}
    \lstinline[basicstyle=\tt]|\usetheme{CambridgeUS}|\\
    \lstinline[basicstyle=\tt]|\usecolortheme{crane}|\\
    \url{https://hartwork.org/beamer-theme-matrix/}\\ \pause
    \lstinline[basicstyle=\tt]|\usefonttheme{structureitalicserif}|\\
    \url{http://deic.uab.es/~iblanes/beamer_gallery/index_by_font.html}\\ 
    \vspace{1cm}
    \inclassFrag{Google \alert{Beamer Theme Matrix} and \alert{Beamer font theme gallery}.}[-1]
\end{frame}

\cprotect\inclassFrame{
\begin{frame}[fragile]{Compile your beamer presentation}\relax
    \only<1>{ \lstinputlisting[basicstyle=\tt\footnotesize, numbers=none]{task1.tex} }
    \only<2>{ \lstinputlisting[basicstyle=\tt\footnotesize, numbers=none]{task2.tex} }
    \only<3>{ \lstinputlisting[basicstyle=\tt\footnotesize, numbers=none]{task3.tex} }
\end{frame}
}

\begin{frame}[fragile]{Colors}
    \skfootnote{\url{en.wikibooks.org/wiki/LaTeX/Colors}}
    \LaTeX provides several standart colors: 
    \textcolor{red}{red}, \textcolor{blue}{blue}, \textcolor{green}{green},\dots\\
    \lstinline[basicstyle=\tt]|\textcolor{red}{text}| \\
    \pause
    There many ways to define new colors, e.~g.
    \lstinline[basicstyle=\tt]|\definecolor{orange}{rgb}{1,.5,0}|\\
    \lstinline[basicstyle=\tt]|\definecolor{orange}{RGB}{255,127,0}|
\end{frame}

\begin{frame}{Colors}
    Beamer automatically loads \alert{xcolor} package\\
    Somehow popular way to define new colors is buy the following rule
    \begin{table}
    \begin{tabular}{0.7\textwidth}\hline
        color           &   rgb formula             &     output  \\\hline
        red!30!blue     &   .3(1,0,0)+.7(0,0,1)     &   \textcolor{red!30!blue}{example} \\
        red!30          &   .3(1,0,0)+.7(1,1,1)     &   \textcolor{red!30!}{example}    \\ 
        red!30!blue!50!green    &   .5(red!30!blue)+.5(0,1,0)   &   \textcolor{red!30!blue!50!green}{example}   
    \end{tabular}          
    \end{table}
\end{frame}

\begin{frame}[t,fragile]{Customazation}
        \lstinline[basicstyle=\tt]|\begin{block}{Block title}|  \\
        \lstinline[basicstyle=\tt]|     Block body|  \\
        \lstinline[basicstyle=\tt]|\end{block}| \\
\end{frame}

\begin{frame}[t,fragile]{Customazation}
        \begin{block}{Block title}
            Block body
        \end{block}
    \pause\vspace{5mm}
    \lstinline[basicstyle=\tt]|\setbeamercolor{block title}{bg=blue!90,fg=white}|
    \lstinline[basicstyle=\tt]|\setbeamercolor{block body}{bg=blue!40!,fg=black}|
    \lstinline[basicstyle=\tt]|\setbeamertemplate{blocks}[rounded][shadow=true]|
    \setbeamercolor{block title}{bg=blue!90,fg=white}
    \setbeamercolor{block body}{bg=blue!40!,fg=black}
    \setbeamertemplate{blocks}[rounded][shadow=true]
    \pause
    \begin{block}{Block title}
          Block body
    \end{block}
\end{frame}

\begin{frame}[fragile]{beamerskoltech.sty}
\skfootnote{github.com/lavton/SkoltechLaTeXtemplates}
    This presentation uses package \\
    \lstinline[basicstyle=\tt]|\usepackage[logo]{beamerskoltech}|\\
    \pause
    This package manages styling and allows to use commands like \
    \lstinline[basicstyle=\tt]|\skfootnote{github.com/lavton/SkoltechLaTeXtemplates}|\\
    \inclassFrag{Upload \texttt{beamerskolthech.sty} from Canvas and compile document with it }[-1]
\end{frame}



\subsection{Inputs}

% \knuthc  knuth the TeXBook
% \lvoc   Lvovsky
% \lamc  lamport latex 
% \slshape different font for footnote
\graphicspath{{sec01/images/s3/}{sec01/code/s3/}}
\lstset{inputpath=sec01/code/s3/}

\begin{frame}[fragile]{\ccol{\footnote}\magicPage}\relax
\cprotect\twocolImg{
    \lstinputlisting[linerange={9-9}]{footnoteMy.tex}
}{footnoteMy}  


\skfootnote{\lmanc{8}[108] \wikiC{https://en.wikibooks.org/wiki/LaTeX/Footnotes_and_Margin_Notes} \overC{https://www.overleaf.com/learn/latex/Footnotes} \knuthc{15}[120]}
\end{frame}

\begin{frame}[fragile]{Horizontal aligment\magicPage}\relax
\cprotect\twocolImg{
    \lstinputlisting[linerange={9-25}]{flushing.tex}
}{flushing}  


\skfootnote{\lmanc{8.12}[63] \lmanc{8.13}[64] \knuthc{14}[112] and note \stExC{https://tex.stackexchange.com/questions/64644/centering-doesnt-seem-to-center-my-text}}
\end{frame}

\begin{frame}[fragile]{Page break\magicPage}\relax

{\centering \Large \ccol{\newpage} \ccol{\pagebreak}\par}

\skfootnote{\lmanc{10}[105] \stExC{https://tex.stackexchange.com/questions/736/pagebreak-vs-newpage}}
     
\end{frame}

\begin{frame}[fragile]{Quotes\magicPage}\relax

\cprotect\twocolImg{
    \lstinputlisting[linerange={10-14}]{quotemy.tex}
}{quotemy}  


\cprotect\skfootnote{\lvoc{III.7.1}[129] \lmanc{8.20}[83]. Also see \verb|quotation| env }
     
\end{frame}

\begin{frame}[fragile]{Verses\magicPage}\relax

\cprotect\twocolImg{
    \lstinputlisting[linerange={11-18}]{versemy.tex}
}{versemy}  

\skfootnote{\lvoc{III.7.3}[130] \lmanc{8.28}[98]}
\end{frame}

\begin{frame}[fragile]{Marginal notes\magicPage}\relax

\cprotect\twocolImg{
    \lstinputlisting[linerange={9-10}]{marginmy.tex}
}{marginmy}  

\skfootnote{\lvoc{IV.10}[194] \lmanc{15.4}[137]}
\end{frame}

 
\cprotect\inclassFrame{
\begin{frame}[fragile]{Try to add a package}\relax
\url{http://hanno-rein.de/downloads/coffee.sty}

and use \verb|\newpage \coffee{2} hello world!|

\incPause
NOW you document is ready!

\end{frame}
}

\section{Programming}
%% It is just an empty TeX file.
%% Write your code here.
\graphicspath{{sec06/images/}{sec06/code/}}
\lstset{inputpath=sec06/code/}

\subsection{Define macros}


\begin{frame}\relax

{\centering\Huge \TeX\ is Turing-complete language

}
\inclassFrag{Who understand, what does it means?}
In basic words, it means, that you can write in \TeX\ and \LaTeX\ any algoritms, that you can write in C++, Java, Python... Moreover, some \TeX syntax is really familiar to functional languages
\skfootnote{\vspace{-3ex}\stExC{https://stackoverflow.com/questions/2968411/ive-heard-that-latex-is-turing-complete-are-there-any-programs-written-in-late} \overC{https://www.overleaf.com/learn/latex/Articles/LaTeX_is_More_Powerful_than_you_Think_-_Computing_the_Fibonacci_Numbers_and_Turing_Completeness} \url{http://sdh33b.blogspot.com/2008/07/icfp-contest-2008.html}}
     
\end{frame}


\begin{frame}{Define macros \tW\magicPage}
     In \TeX\ you can define new macros via \ccol\def.

\twocolImg{
\lstinputlisting[linerange={8-9, 14-16}]{commandwith.tex}
}{commandwith}
     
    Use \ccol\global\ prefix to define macros not just inside ``group''. 
    
    Use \ccol\long\ prefix to define macros that can have multiple paragraphs as an argument.
    
    \skfootnote{\tugC{https://www.tug.org/utilities/plain/cseq.html\#def-rp} \knuthc{20}[209]}
\end{frame}

\begin{frame}{Define with pattern matching\tw\magicPage}\relax

The syntax with writing each argument seems to be an over-use. But it is needed because of \textit{pattern matching}
  
\twocolImg{
\lstinputlisting[linerange={11-13}]{commandpattern.tex}
}{commandpattern}   

\end{frame}

\subsection{Conditions}

\begin{frame}[fragile]{Compare strings (macros)\tW\magicPage}\relax

\twocolImg{
\lstinputlisting[linerange={12-21}]{ifmy.tex}
}{ifmy}   

    {\csk \verb|\ifx\<first>\<second>  <code1>  [\else  <code2>]  \fi|}
    
\end{frame}

\begin{frame}[fragile]{Compare numbers\tW\magicPage}\relax

\twocolImg{
\lstinputlisting[linerange={12-20}]{ifnummy.tex}
}{ifnummy}   

    {\csk \verb|\ifnum<first><operator><second>  <code1>  [\else  <code2>]  \fi|}. Only ``='', ``>'' or ``<'' are allowed.
    
    Use {\csk \verb|\ifcase|} to check different stuff.
    
    Also use {\csk \verb|\ifodd|} to check if num is odd or even
    
\end{frame}

\begin{frame}[fragile]{Check modes\tW\magicPage}\relax

\twocolImg{
\lstinputlisting[linerange={12-22}]{ifmodemy.tex}
}{ifmodemy}   
\begin{itemize}
    \item {\csk\verb|\ifmmode|} to check if in mathematical mode
    \item {\csk\verb|\ifvmode|} to check if in vertical mode
    \item {\csk\verb|\ifhmode|} to check if in horizontal mode
    \item {\csk\verb|\ifinner|} to check if \TeX\ is in internal vertical mode, or restricted horizontal mode, or (nondisplay) mathmode
     
\end{itemize}

\skfootnote{other ``if'''s can be found in \knuthc{20}[221]}
\end{frame}


\begin{frame}{Compare in \LaTeX\lW\magicPage}\relax
\twocolImg{
\lstinputlisting[linerange={7-7, 12-16}]{ifstring.tex}
}{ifstring}  
      \ncol\usepackage{xstring}
      
      also see \ncol\usepackage{ifthen}
      
      \hrule 
      
      you can check if you are in \XeLaTeX by \ncol\usepackage{ifxetex}
      
      \skfootnote{read about the xstring package! It has lots of things: strings, substrings,.. what to do with them,.. replace strings etc\\ \stExC{https://tex.stackexchange.com/questions/47576/combining-ifxetex-and-ifluatex-with-the-logical-or-operation}}
\end{frame}

\subsection{Loops and recursion}

\begin{frame}[fragile]{Loop \tW\magicPage}\relax
\twocolImg{
\lstinputlisting[linerange={12-17}]{loopmy.tex}
}{loopmy}

\ccol\loop\ for start loop, <code> inside, then {\csk\verb|\if<..>|}-family, another bunch of <code>, ended with \ccol\repeat.

     \skfootnote{\knuthc{20}[228]}
\end{frame}

\begin{frame}[fragile]{Reqursion \tW\magicPage}\relax
\twocolImg{
\lstinputlisting[linerange={12-13}]{reqmy.tex}
}{reqmy}

\end{frame}


\begin{frame}[fragile]{For loop\lW\magicPage}\relax
\twocolImg{
\lstinputlisting[linerange={8-8, 13-16}]{forloopmy.tex}
}{forloopmy}
\ncol \usepackage{forloop}

\twocolImg{
\lstinputlisting[linerange={8-8, 13-15}]{foreachmy.tex}
}{foreachmy}

\ncol \usepackage{pgffor}, part of pgf, part of TikZ

\skfootnote{\stExC{https://stackoverflow.com/questions/2561791/iteration-in-latex}}
\end{frame}

\subsection{Related things to macros creation}

\begin{frame}[fragile]{``let'' command\magicPage}\relax

\twocolImg{
\lstinputlisting[linerange={11-21}]{letmy.tex}
}{letmy}


The statement ``\ccol\let\verb|\a=\b|'' gives \string\a\ the current meaning of \string\b. If \string\b\ changes after the assignment is made, \string\a\ does not change.

\skfootnote{\knuthc{20}[217] \tugc{https://www.tug.org/utilities/plain/cseq.html\#let-rp}}
     
\end{frame}

\begin{frame}[fragile]{Usecase with \ccol\let: ``decorator''\magicPage}\relax

Imagine: you have some \string\command used inside the document multiple times. You want to \textit{add some addition behaviour} to the command -- decorate (or ``wrap'', or ``redefine with the use of itself''). You can do it with \ccol\let:

\begin{lstlisting}
\let\oldCommand=\command
\def\command#1{<some code>\oldCommand}
\end{lstlisting}

And the same for enviruments using \ncol\usepackage{etoolbox} or \ccol{\g@addto@macro}



\skfootnote{\stExC{https://tex.stackexchange.com/questions/47351/can-i-redefine-a-command-to-contain-itself} \stExC{https://tex.stackexchange.com/questions/467435/environment-decorators}}
     
\end{frame}


\subsection{Programming examples}

\begin{frame}{99 Bottles of Beer\magicPage}\relax

\twocolImg{
\lstinputlisting[linerange={9-26}]{bottles.tex}
}{bottles}

\skfootnote{\url{https://rosettacode.org/wiki/99_Bottles_of_Beer#LaTeX}}
     
\end{frame}

\begin{frame}{not-AND logical gate\magicPage}\relax

\twocolImg{
\lstinputlisting[linerange={11-29}]{nand.tex}
}{nand}

\skfootnote{\url{https://rosettacode.org/wiki/99_Bottles_of_Beer#LaTeX}}
\end{frame}

\begin{frame}{Split words\magicPage}\relax
\twocolImg{
\lstinputlisting[linerange={11-31}, basicstyle=\tiny]{splitmy.tex}
}{splitmy}

     \skfootnote{\stExC{https://tex.stackexchange.com/questions/12810/how-do-i-split-a-string/12811}}
\end{frame}

\subsection{Relatate stuff}
\graphicspath{{sec08/images/}{sec08/code/}}
\lstset{inputpath=sec08/code/}


\begin{frame}{Filesystem writing\magicPage}
\twocolImg{
\lstinputlisting[linerange={11-17}]{writedoc.tex}
}{writedoc}

\footnotesize
The output in pdf is the result of listing the {\csk outfile.txt}
\begin{itemize}
    \item \ccol\newrite\ to get a free register
    \item \ccol\openout\ to open filename for output
    \item \ccol\write<register>\ to actually write 
    \item \ccol\noexpand\ to stop expanding (also see \ccol\expandafter)
    \item \ccol\closeout\ to close the file 
\end{itemize}

You also can use numbers. Like \ccol\write4

\end{frame}

\begin{frame}[fragile]{Filesystem reading\magicPage}
\twocolImg{
\lstinputlisting[linerange={18-31}, basicstyle=\tt\tiny]{readdoc.tex}
}{readdoc}

\footnotesize
\begin{itemize}
    \item \ccol\newread\ to get a free register
    \item \ccol\openin\ to open filename for input
    \item \ccol\read<register>\ \ccol{to}\string\<newvariable> to actually read
    \item \verb"\ifeof" to check if file is still have lines
    \item \ccol\closein\ to close the file 
\end{itemize}

You also can use numbers. Like \ccol\read4

\end{frame}


\begin{frame}[fragile]{How BiBLaTeX works? Proof-of-concept\magicPage}
\twocolImg{
\lstinputlisting[linerange={8-22, 28-31}, basicstyle=\tt\tiny]{PoCIO.tex}
}{PoCIO}

\footnotesize
\begin{enumerate}
    \item write info into a file 
    \item use an external command to do something with the file 
    \item read content from a file in a different place
     
\end{enumerate}

\end{frame}

\begin{frame}[fragile]{Use command line\magicPage}\relax
    \twocolImg{
\lstinputlisting[linerange={12-14}]{commandline.tex}
}{commandline}
\small

Use 
\begin{itemize}
    \item \ccol\write18\ to call the command line 
    \item \ccol\immidiate\ to run it as it reached (otherwise only when \TeX\ will ``print'' the page)
    \item use \verb|--enable-write18 -interaction=nonstopmode| keys for run offline
    \item the commands with internet connection will not work at Papeeria
     
\end{itemize}
\skfootnote{\stExC{https://tex.stackexchange.com/questions/20444/what-are-immediate-write18-and-how-does-one-use-them}. Also look at \ccol{filecontents} env.}
\end{frame}

\begin{frame}{Command names manipulation\magicPage}\relax
    \twocolImg{
\lstinputlisting[linerange={12-13}]{csnamemy.tex}
}{csnamemy}

\begin{itemize}
    \item \ccol\csname\ --- \ccol\endcsname\ to ``compile'' command from name 
    \item \ccol\string\ to show the name 
     
\end{itemize}

\end{frame}

\begin{frame}{Catcodes\magicPage}\relax

\twocolImg{
\lstinputlisting[linerange={12-18}]{catcodesmy.tex}
}{catcodesmy}

\ccol\catcode\ shows what symbol will be responsible for the group, what for comment etc.

    \skfootnote{\tugc{https://www.tug.org/utilities/plain/cseq.html#catcode-rp} \knuthc{app. B}[354]}
\end{frame}

\begin{frame}{Lua\TeX\ Proof-of-concept\magicPage}\relax
\twocolImg{
\lstinputlisting[linerange={8-8, 10-13}]{luamy.tex}
}{luamy}
     
\skfootnote{\stExC{https://tex.stackexchange.com/questions/70/what-is-a-simple-example-of-something-you-can-do-with-luatex}}
\end{frame}

% %write/read, expand,...

\section{Debugging}
%% It is just an empty TeX file.
%% Write your code here.
\graphicspath{{sec05/images/}{sec05/code/}}
\lstset{inputpath=sec05/code/}

\newcommand{\logshowI}[3]{ %%% add two columns: (1) some latex-code (2) the log file text produced by the code
    \begin{columns}
    \begin{column}{0.45\textwidth}
        %  \lstinputlisting[linerange={#2}]{#1.tex}
          \foreach \range in {#2} {%
                \expandafter\StrBefore{\range}{-}[\start]%  % Extract the start value
                \expandafter\StrBehind{\range}{-}[\eend]%   % Extract the end value
                %
                \expandafter\myminted\expandafter{sec05/code/#1.tex}{\expandafter\start}{\expandafter\eend}%
            }
    \end{column}
    
    \begin{column}{0.45\textwidth}
         \lstinputlisting[language={},linerange={#3}, keywordstyle=\color{black}, identifierstyle=\color{black},]{#1.logp}
    \end{column}
    \end{columns}
}

\subsection{Show-family}

\begin{frame}[fragile]{Expand macros\magicPage}\relax

    \logshowI{showmy}{7-10,38-41}{1-18}
    
    \ccol{\show<macros>} 
     \skfootnote{\tugC{https://www.tug.org/utilities/plain/cseq.html\#show-rp} \knuthc{3}[21] \stExC{https://tex.stackexchange.com/questions/364942/show-family-debugging-commands}}
\end{frame}

\begin{frame}[fragile]{Expand macros\magicPage}{\ccol{show} with not macros}\relax
    
    \logshowI{showmy}{12-15,43-44}{21-26}
\end{frame}

\begin{frame}[fragile]{Show length and counts\tW\magicPage}\relax

    \logshowI{showthemy}{12-20,39-45}{1-14}
    \ccol{\showthe<var>}
    % TODO: \showthe\font
    \skfootnote{\url{https://texfaq.org/FAQ-printvar} \tugC{https://www.tug.org/utilities/plain/cseq.html\#showthe-rp} \knuthc{15}[132]}
\end{frame}


\begin{frame}[fragile]{Show boxes\tW\magicPage}\relax

    \logshowI{showboxmy}{31-39}{1-14}
    \ccol{\showbox<box>}
    
    \skfootnote{\url{https://texfaq.org/FAQ-printvar} \tugC{https://www.tug.org/utilities/plain/cseq.html\#showbox-rp} \tugC{https://www.tug.org/utilities/plain/cseq.html\#showboxbreadth-rp} \tugC{https://www.tug.org/utilities/plain/cseq.html\#showboxdepth-rp} \knuthc{11}[77]\\ \ccol{\showboxbreadth} and \ccol{\showboxdepth} provides how much info to show}
\end{frame}

% \begin{frame}[fragile]{Show stack of modes\tW\magicPage}\relax

% \logshow{showlistsmy}{12-17}{1-18}
% \ccol{\showlists}
% \skfootnote{\vspace{-3ex}\tugC{https://www.tug.org/utilities/plain/cseq.html\#showlists-rp}  \knuthc{13}[99]\\``it causes TEX to display the lists that are being worked
% on, in the current mode and in all enclosing modes where the work has been suspended:''}
% \end{frame}

\begin{frame}{Show-family list}\relax
    \small
    \begin{description}
        \item[\ccol{\show}] log macros insides
        \item[\ccol\showthe] log length or counter value
        \item[\ccol\showbox] log box insides 
        \begin{description}
            \item[\ccol\showboxdepth] the value of the deepest level of box nesting
            \item[\ccol\showboxbreadth] the maximum number of items shown per level 
        \end{description}
        \item[\ccol\showlists] writes the content of partial box lists in all of the 4 non-math TeX modes
        \item[\ccol\showhyphens\{W\}] displays the hyphenation of W on the terminal/log according to the hyphenation rules.
    \end{description}
     \skfootnote{\stExC{https://tex.stackexchange.com/questions/364942/show-family-debugging-commands}}
\end{frame}

\subsection{Tracing-family}

\begin{frame}[fragile]{Trace modes and commands\tW\magicPage}\relax
    \logshowI{tracingcommandsmy}{11-11,16-18}{1-10}
    \ccol{\tracingcommands=1}
    
    \skfootnote{\tugC{https://www.tug.org/utilities/plain/cseq.html\#tracingcommands-rp} \knuthc{13}[99]}
\end{frame}

\begin{frame}[fragile]{Trace macros (recursively)\tW\magicPage}\relax
    \logshowI{tracingmacrosmy}{7-10,15-20}{1-15}
    \ccol{\tracingmacros=1}
    
    \skfootnote{\tugC{https://www.tug.org/utilities/plain/cseq.html\#tracingmacros-rp} \knuthc{20}[216]}
\end{frame}


\begin{frame}{Tracing-family list}\relax
    \footnotesize
    
    \begin{description}
        \item[\ccol\tracingcommands] if positive, writes commands to the .log file
        \item[\ccol\tracinglostchars] if positive, writes characters not in the current font to the .log file
        \item[\ccol\tracingmacros] if positive, writes to the .log file when expanding macros and arguments
        \item[\ccol\tracingonline] if positive, writes diagnostic output to the terminal as well as to the .log file
        \item[\ccol\tracingoutput] if positive, writes contents of shipped out boxes to the .log file
        \item[\ccol\tracingpages] if positive, writes the page-cost calculations to the .log file
        \item[\ccol\tracingparagraphs] if positive, writes a summary of the line-breaking calculations to the .log file
        \item[\ccol\tracingrestores] if positive, writes save-stack details to the .log file
        \item[\ccol\tracingstats] if positive, writes memory usage statistics to the .log file
        \item[\ccol\tracingall] turns on every possible mode of interaction
    \end{description}
     \skfootnote{\stExC{https://tex.stackexchange.com/questions/60491/latex-tracing-commands-list}}
\end{frame}

\subsection{Other debugging ways}

\begin{frame}[fragile]{Message to log file\magicPage}\relax
    \logshowI{messagemy}{7-7,14-15,17-18,31-32,15-20}{1-2}
    \ccol{\message\{<msg>\}} -- \TeX-command, \ccol{\typeout\{<msg>\}} -- \LaTeX-command
    
    \skfootnote{\tugC{https://www.tug.org/utilities/plain/cseq.html\#message-rp} \knuthc{20}[228] \lmanc{27.2}[228]}
\end{frame}

% 

\graphicspath{{code_ex/}}
\lstset{inputpath=code_ex/}
\cprotect\inclassframe{
\begin{frame}{\exFrame{Implement for loop}}{lvl 1}

     Implement a function that, given the syntax in the code on the left, would output the result as in the image on the right.
    \twocolImg{
    \inputminted[firstline=24, lastline=27]{latex}{code_ex/exForLvl1.tex}
    }{exForLvl1.pdf}
    
    Tips: 
    \begin{itemize}
        \item You need a counter
        \item Use \ccol\foreach
        \item if you increment counter, but it remains the same, try increment in with \ccol\global
    \end{itemize}
\end{frame}
\begin{frame}[fragile]{\exFrame{Implement for loop}}{lvl 2}
     Implement a function that, given the syntax in the code on the left, would output the result as in the image on the right. You can use \ccol\foreach, but not any of addition packages
    \twocolImg{
    \inputminted[firstline=24, lastline=28]{latex}{code_ex/exForLvl2.tex}
    }{exForLvl2.pdf}
    
    Tips: 
    \begin{itemize}
        \item The problem will be how to to split \{2-5\}\{Caption 1\}
        \item You may need a \ccol\expandafter
        \item Create a helper function that would do 1 line and play around with how you can pass parameters to it
    \end{itemize}
\end{frame}
\begin{frame}[fragile]{\exFrame{Implement for loop}}{lvl 2. Hardcore mode!}
     Implement a function that, given the syntax in the code on the left, would output the result as in the image on the right. You can't use \ccol\foreach\ or any of addition packages! Try use \TeX\ power
    \twocolImg{
    \inputminted[firstline=31, lastline=35]{latex}{code_ex/exForLvl2.5.tex}
    }{exForLvl2.pdf}
    
    Tips: 
    \begin{itemize}
        \item You need to use pattern matching tehnique
        \item You need to use reqursion
        \item One of your helper macros may have the following declaration: \verb|\def\<funName>#1#2, #3\relax|
    \end{itemize}
\end{frame}
}
\cprotect\inclassframe{
\begin{frame}[fragile]{\exFrame{Implement for loop}}{lvl 1: more tips}\relax
\pause
\begin{itemize}
    \item Firstly try to make it without counter \pause
    \item The \ccol\foreach\ command was in the slide \ref{sl:for} \pause
    \item There the format like ``\verb|apples,burgers,cake|'', you have ``\verb|{Caption 1}, {Caption, 2}|''. But braces just indicates atomicy (slide \ref{sl:atom})! \pause
    \item So, ``\verb|\foreach \n in {apples,burgers,cake}|'' is the same as ``\verb|\foreach \n in {{Caption 1}, {Caption, 2}}|''\pause
    \item The only difference if you just pass ``\verb|\foreach \n in #1|'' it will be expanded as ``\verb|\foreach \n in {Caption 1}, {Caption, 2}|'', not ``\verb|\foreach \n in {{Caption 1}, {Caption, 2}}|''. Guess what to do!\pause
    \item After stuff without counter works, try to add a counter. How to make counter see at previous lecture. Don't forget to use \ccol\global\ if you try to add 1 to it at it is not working
\end{itemize}
     
\end{frame}
}

\progressend


\begin{frame}\frametitle{What we have learned today?}\relax
    \tableofcontents
\end{frame}

\begin{frame}[allowframebreaks]{references}
color from the footnotes corresponds to references' color.
    \begin{itemize}
        \item \knuthc{Knuth ``The \TeX Book''}
        \item \lvoc{L'vovsky ``Nabor i verstka v sisteme \LaTeX''}
        \item \lamc{Lamport. ``\LaTeX. A Document Preparation System, User’s Guide and Reference Manual''}
        \item \lmanc{``\LaTeX 2e: An unofficial reference manual''} also at website \url{https://latexref.xyz/}
        \item \stExC{https://tex.stackexchange.com/questions} : \url{https://tex.stackexchange.com/questions}
        \item \wikiC{https://en.wikibooks.org/wiki/LaTeX} : \url{https://en.wikibooks.org/wiki/LaTeX}
        \item \overC{https://www.overleaf.com/learn/latex} : \url{https://www.overleaf.com/learn/latex}
        \item \tugC{https://www.tug.org/utilities/plain/cseq.html} : \url{https://www.tug.org/utilities/plain/cseq.html}
    \end{itemize}
\end{frame}

\begin{frame}{Distribution}\relax
\begin{itemize}
     \item the pdf-version of the presentation and all printed materials can be distributed under license Creative Commons Attribution-ShareAlike 4.0 \url{https://creativecommons.org/licenses/by-sa/4.0/}
     \item The source code of the presentation is available on {\csk\url{https://github.com/Lavton/latexLectures}} and can be distributed under the MIT license \url{https://en.wikipedia.org/wiki/MIT_License\#License_terms}
\end{itemize}
     
\end{frame}
\end{document}